\documentclass[11pt]{article}

\usepackage[T1]{fontenc}
\usepackage[utf8]{inputenc}
\usepackage{amsmath,amssymb,amsthm,mathtools}
\usepackage{enumitem}
\usepackage{microtype}
\usepackage[colorlinks=true,linkcolor=blue,citecolor=blue,urlcolor=blue]{hyperref}

\hypersetup{
  pdftitle={A dense sequence with distinct products of consecutive blocks},
  pdfauthor={Fuzzy Dynamics}
}

\title{A dense sequence with distinct products of consecutive blocks}
\author{Fuzzy Dynamics}
\date{\today}

\newtheorem{theorem}{Theorem}[section]
\newtheorem{proposition}[theorem]{Proposition}
\newtheorem{lemma}[theorem]{Lemma}
\newtheorem{corollary}[theorem]{Corollary}
\newtheorem{external}[theorem]{External input}
\theoremstyle{definition}
\newtheorem{definition}[theorem]{Definition}
\newtheorem{remark}[theorem]{Remark}

\newcommand{\N}{\mathbb N}
\newcommand{\Z}{\mathbb Z}
\newcommand{\Q}{\mathbb Q}
\newcommand{\A}{\mathbb A}
\newcommand{\eps}{\varepsilon}
\newcommand{\pplus}{P^+}
\newcommand{\floor}[1]{\left\lfloor #1\right\rfloor}
\newcommand{\ceil}[1]{\left\lceil #1\right\rceil}
\newcommand{\abs}[1]{\left|#1\right|}
\newcommand{\norm}[1]{\left\|#1\right\|}
\newcommand{\1}{\mathbf 1}
\DeclareMathOperator{\rad}{rad}
\DeclareMathOperator{\height}{ht}

\begin{document}
\maketitle

\begin{abstract}
We construct a set \(A\subset\N\) of natural density one whose increasing
enumeration \(d_1<d_2<\cdots\) has the property that the products over finite
consecutive blocks are all distinct.  The set is produced by a greedy
algorithm, modified by a deterministic deletion set of density zero.  The role
of this deletion set is to convert the semiprimes supplied by the
Matom\"aki--Ter\"av\"ainen theorem into private barriers: short obstructions
carrying a globally unique large prime.  Once a few density-zero exceptional
families are discarded, the barriers confine each surviving collision to a
polylogarithmic window, and the resulting product identities carry many
large-prime congruences of shifted-CRT type.  A low-degree auxiliary
polynomial interpolating the local data is then forced to vanish at the global
collision point.  Its nonlinear components are counted by the
Castryck--Cluckers--Dittmann--Nguyen affine-curve bound, and its line
components directly from the curvature of \(q^\theta\).  The outcome is
\(o(X)\) surviving collision endpoints in each dyadic shell \((X,2X]\), and
hence a density-one set with distinct consecutive-block products.
\end{abstract}

\medskip
\noindent\emph{2020 Mathematics Subject Classification.}
Primary 11B75; Secondary 11N25, 11N35, 11D45.

\smallskip
\noindent\emph{Key words and phrases.}
Distinct products, consecutive blocks, natural density, greedy construction,
smooth numbers, determinant method, Erd\H{o}s problems.

\section{Introduction}

An increasing sequence of positive integers assigns to each finite interval of
indices the product of the corresponding terms, and one may ask how rigidly
those block products pin down the interval.  For an increasing sequence
\(d_1<d_2<\cdots\), write
\[
  P_d(u,v)=\prod_{i=u}^v d_i \qquad (1\le u\le v).
\]
Say that the sequence has \emph{distinct consecutive-block products} if
\[
  P_d(u,v)=P_d(u',v') \quad\Longrightarrow\quad (u,v)=(u',v')
\]
for all finite index intervals, so that no two different blocks share a
product.  Such rigidity is easy to arrange in a thin sequence; the question is
how dense one can make it.  Erd\H{o}s asked whether it is compatible with
natural density one.  This is Problem 421 in the Erd\H{o}s problem collection
\cite{Bloom421}, going back to the problem list of Erd\H{o}s and Graham
\cite{ErGr80}.

\begin{theorem}\label{thm:main}
There exists a strictly increasing sequence
\[
  1\le d_1<d_2<\cdots
\]
whose underlying set has natural density \(1\), and for which the products
\(\prod_{i=u}^v d_i\) are all distinct as \((u,v)\) ranges over finite
consecutive intervals.
\end{theorem}

Our argument is self-contained except for the external results collected in
Section~\ref{sec:inputs}.  The sequence is built greedily, with one
modification.  We fix in advance a deterministic set \(D\) of density zero,
reject every integer of \(D\) outright, accept \(2\), and accept each later
integer outside \(D\) precisely when doing so creates no new coincidence among
consecutive-block products.  Uniqueness is then automatic from the greedy
rule, and the whole difficulty lies in showing that the rejected integers --
those forced by \(D\) together with those the greedy test discards -- still
form a set of density zero.

To locate that difficulty, suppose an integer \(n\notin D\) is rejected.  Then
adjoining \(n\) produces a collision, and a short analysis
(Lemma~\ref{lem:witness-split}) shows that the collision is always an
\emph{old-new witness}
\[
  P=nS,
\]
in which \(P\) is a block product already present and \(S\) is a product of
the most recent accepted terms, possibly empty.  When \(S=1\), the integer
\(n\) is itself an old block product, and such \(n\) are sparse for elementary
reasons.  The genuine obstruction is the case \(S>1\): here \(P\) may be long
and may sit far to the left of the suffix, so no direct count is available.
The set \(D\) is designed precisely to force almost every such witness into a
short window, where it can be counted.

This window is produced by what we call a \emph{barrier system}.  In most blocks
of polylogarithmic length we single out one semiprime \(b=p_1p_2\) with
\((\log b)^{1.08}<p_1\le(\log b)^{1.1}\) whose large prime factor \(p_2\) is
globally unique among our choices, and we delete every other multiple of
\(p_2\) up to a long but finite horizon.

Two properties make this effective.  A selected barrier survives the greedy
scan, because any new product through it inherits the private prime \(p_2\)
that no earlier product can carry; hence the barrier is accepted.  And once a
barrier is in place, no surviving collision can have its old block cross it,
unless the endpoint has already passed the barrier's horizon.  The integers
beyond their horizon make up a smooth family of density zero, removed
separately.

With the barriers in place, every nonempty collision left in a shell
\((X,2X]\) has a short, tightly controlled witness
\[
  \prod_{i=1}^{\ell}(x-a_i)=\prod_{j=1}^{b}(q+c_j),
  \qquad 0\le a_i,\, c_j \le (\log X)^{O(1)},
\]
with \(x<q<n\le 2X\), \(q\) a previously selected barrier, and
\[
  x=q^{b/\ell}+O((\log X)^B),\qquad
  (\log X)^{-B}\le 1-\frac b\ell \ll \frac{\log\log X}{\log X}.
\]
All the shifts \(a_i,c_j\) are now bounded by a fixed power of \(\log X\), and
the argument becomes arithmetic.  A large prime dividing a factor on the left
must reappear among the terminal factors on the right, and matching these
yields many congruences
\[
  p_i\mid x-a_i,\qquad p_i\mid q-x+a_i+c_i,
\]
with distinct \(p_i>\exp(\log X/(\log\log X)^2)\).  A low-degree auxiliary
polynomial that vanishes on the local points \((a_i,a_i+c_i)\) is therefore
forced to vanish at the global point \((x,x-q)\) as well.  Its nonlinear
components are counted by the affine-curve determinant-method theorem of
Castryck--Cluckers--Dittmann--Nguyen, while its line components are handled
directly, through the strict curvature of \(q^\theta\) for \(\theta<1\).

The paper is organized as follows.  Section~\ref{sec:inputs} fixes the
external inputs and notation.  Section~\ref{sec:greedy} records the formal
greedy reduction and disposes of empty-suffix collisions.
Section~\ref{sec:barriers} constructs the barrier system.
Sections~\ref{sec:localization} and~\ref{sec:terminal} turn the barrier
information into localized product identities and then into terminal
large-prime support.  Section~\ref{sec:slab} proves the shifted-CRT counting
proposition.  Section~\ref{sec:nonempty} applies it to all remaining nonempty
rejections, and Section~\ref{sec:main-proof} completes the density argument.

\section{External Inputs and Notation}\label{sec:inputs}

All logarithms are natural, except \(\log_2\).  Constants implicit in
\(O(\cdot)\), \(o(\cdot)\), and \(\ll\) may depend on the fixed constants
chosen in the proof, but not on the dyadic parameter \(X\).  A \emph{dyadic
shell} means \((X,2X]\), with \(X\) a sufficiently large power of two.
For an integer \(m\ge2\), let \(\pplus(m)\) be its largest prime factor.  Let
\[
  \Psi(X,Y)=\#\{m\le X:\pplus(m)\le Y\}.
\]

\begin{external}[Matom\"aki--Ter\"av\"ainen \(E_2\) theorem \cite{MT}]\label{inp:MT}
There are constants \(c_0,\delta>0\) such that, for all sufficiently large
\(X\), all but \(O(X/(\log X)^\delta)\) integer starts \(2\le y\le X\) satisfy
\[
\#\left\{p_1p_2\in (y,y+(\log y)^{2.1}]:
\begin{array}{c}
p_1,p_2\ {\rm prime},\\
(\log y)^{1.09}<p_1\le(\log y)^{1.1}
\end{array}\right\}
\ge c_0(\log y)^{1.1}.
\]
\end{external}

\begin{external}[de Bruijn--Dickman smooth-number estimate
\cite{deBruijn,Hildebrand}]\label{inp:dickman}
Uniformly for \(u=u(N)\to\infty\) with \(u=O(\sqrt{\log\log N})\),
\[
  \Psi(N,N^{1/u})\le
  N\exp\big(-(1+o(1))u\log u\big).
\]
\end{external}

\begin{external}[CCDN affine-curve bound \cite{CCDN}]\label{inp:CCDN}
For every fixed ambient dimension \(n>1\), there is a constant \(C_n\) such
that, for every \(B\ge1\), every integral affine curve \(V\subset\A^n_{\Q}\) of
positive degree \(d\) satisfies
\[
  \#\{z\in V(\Z):\height(z)\le B\}
  \le C_n d^3 B^{1/d}(\log B+d).
\]
Here \(\height(z)=\max_i |z_i|\) for \(z=(z_1,\ldots,z_n)\in\Z^n\).
The constant \(C_n\) is independent of \(d\), \(B\), and the defining
equations of \(V\); the displayed factor is the only degree-dependence used
below.
\end{external}

\begin{external}[Mignotte factor inequality \cite{Mignotte}]\label{inp:mignotte}
If \(f,g\in\Z[z]\) are nonzero, \(\deg g\le d\), and \(g\) is a multiple of
\(f\), then
\[
  \norm{g}_2\ge 2^{-d}\norm{f}_\infty .
\]
\end{external}

We also use the following standard elementary facts: Bertrand's postulate,
Chebyshev upper and lower prime-counting bounds, Gauss's lemma for primitive
factors over \(\Q\) and \(\Z\), the usual compatibility of primitive
homogenization and dehomogenization for line factors, Galois conjugacy of
absolute factors of a squarefree \(\Q\)-irreducible plane curve, and Bezout's
\(d_1d_2\) bound for two coprime plane curves of degrees \(d_1,d_2\).

\section{The Modified Greedy Reduction}\label{sec:greedy}

It is convenient to isolate the purely formal consequences of the greedy rule
from the number-theoretic construction of \(D\), which is deferred to
Section~\ref{sec:barriers}.  Granting that \(D\) has density zero, the only
rejections that will demand real work are the nonempty old-new collisions
introduced below.

\begin{definition}[Modified greedy sequence]\label{def:modified-greedy}
Fix a deterministic set \(D\subset\N\) with \(1\in D\), \(2\notin D\), and
\(\abs{D\cap[1,N]}=o(N)\).  The modified greedy construction rejects every
\(n\in D\), accepts \(2\), and for each \(n\ge3\) with \(n\notin D\), accepts
\(n\) if and only if adjoining \(n\) to the accepted prefix creates no equality
between two distinct finite consecutive-block products.
\end{definition}

\begin{lemma}[Greedy uniqueness]\label{lem:greedy-unique}
Every finite accepted prefix produced by Definition~\ref{def:modified-greedy}
has distinct consecutive-block products.
\end{lemma}

\begin{proof}
The initial accepted prefix is \(\{2\}\).  A rejection leaves the accepted
prefix unchanged.  An acceptance occurs only after the enlarged prefix has
been tested to have no equality between two distinct consecutive-block
products.  Induction over the scan proves the claim.
\end{proof}

\begin{lemma}[Old-new witness split]\label{lem:witness-split}
Let \(n\ge3\) be rejected by the modified greedy construction and suppose
\(n\notin D\).  If \(d_1<\cdots<d_m\) is the accepted prefix before testing
\(n\), then there is an old consecutive accepted block product \(P\) and a
suffix product \(S\) of the old prefix, possibly \(S=1\), such that
\[
  P=nS.
\]
\end{lemma}

\begin{proof}
The old prefix has distinct block products by Lemma~\ref{lem:greedy-unique}.
Since \(n\notin D\), rejection means that adding \(n\) creates a collision.
Old-old collisions are impossible.  Two new block products both contain \(n\);
after cancelling \(n\), one obtains equality of two final suffix products of
the old prefix.  Distinct final suffixes are strictly ordered by inclusion,
and extending a suffix multiplies its product by an accepted integer at least
\(2\).  Hence two new products cannot collide.  The collision is therefore
between an old block and a new block, and every new block is \(n\) times an
old suffix product.
\end{proof}

\begin{definition}
A nonforced rejection \(n\notin D\) is an \emph{empty-suffix rejection} if it
has some witness in Lemma~\ref{lem:witness-split} with \(S=1\).  Otherwise we
choose, once and for all, the lexicographically first old-new witness in
accepted-index coordinates, first by old-block endpoints and then by suffix
start, and call \(n\) a \emph{nonempty rejection}.  For nonempty rejections we
cancel the maximal common contiguous overlap between the old block and the
final suffix.  The old block is not contained in that suffix, and if the
cancellation exhausted the suffix then \(n\) would itself be an old block
product.  Thus the remaining suffix is nonempty, and the old block lies
strictly before it.
\end{definition}

\begin{lemma}[Sparse old products]\label{lem:block-products-sparse}
For any increasing positive-integer sequence \(d_1<d_2<\cdots\) with
\(d_1\ge2\), the number of values at most \(N\) which occur as
\(P_d(u,v)\) with \(v-u+1\ge2\) is \(o(N)\).
\end{lemma}

\begin{proof}
Since \(d_i\ge i+1\), a length-\(\ell\) block starting at index \(i\) has
product at least \((i+1)^\ell\).  For fixed \(\ell\ge2\), at most
\(N^{1/\ell}\) length-\(\ell\) block products can be \(\le N\).  No length
\(\ell>\log_2 N\) contributes.  Thus the number of such values is at most
\[
  \sum_{2\le \ell\le \log_2 N} N^{1/\ell}
  =O(N^{1/2}\log N)=o(N).
\]
\end{proof}

\begin{corollary}\label{cor:empty-suffix}
The empty-suffix rejections of the modified greedy construction have density
zero.
\end{corollary}

\begin{proof}
An empty-suffix rejection \(n\) is an old block product.  The old block has
length at least two, since length one would make \(n\) an already accepted
old term.  Apply Lemma~\ref{lem:block-products-sparse} to the accepted
sequence.
\end{proof}

\begin{lemma}[Elementary smooth decay]\label{lem:smooth-decay}
Let \(E>0\), \(\rho>1\), and \(M>0\) be fixed.  Then
\[
  \Psi(X,(\log X)^E)=o(X)
\]
and, with \(Y_X=\exp(\log X/(\log\log X)^\rho)\),
\[
  \Psi(2X,Y_X)=o\!\left(\frac{X}{(\log X)^M}\right).
\]
\end{lemma}

\begin{proof}
For \(1/2\le\sigma<1\), Rankin's trick gives
\[
  \Psi(X,Y)\le X^\sigma\prod_{p\le Y}(1-p^{-\sigma})^{-1}.
\]
Write \(\sigma=1-\eta\), \(0<\eta\le1/2\).  Chebyshev's bound and partial
summation give
\[
\log\prod_{p\le Y}(1-p^{-\sigma})^{-1}
 \ll \log\log(3Y)+\frac{Y^\eta}{\eta\log Y}.
\]
This rests on the estimate
\[
\sum_{p\le Y}p^{-1+\eta}
\le \pi(Y)Y^{-1+\eta}
 +(1-\eta)\int_2^Y \pi(t)t^{-2+\eta}\,dt
\ll \frac{Y^\eta}{\log Y}
 + \int_2^Y \frac{t^{-1+\eta}}{\log t}\,dt,
\]
and splitting the last integral at \(\log t=1/\eta\), which gives
\(O(\log\log(3Y))\) below the split and \(O(Y^\eta/(\eta\log Y))\) above it.

For \(Y=(\log X)^E\), choose fixed \(\sigma\in(0,1)\) with
\(E(1-\sigma)<1\), obtaining \(X^\sigma\exp(o(\log X))=o(X)\).
For \(Y_X=\exp(L/u)\), \(L=\log X\), \(u=(\log\log X)^\rho\), take
\(\eta=(\log u)/(2\log Y_X)\).  Then \(Y_X^\eta=\sqrt u\) and
\(\eta L=\tfrac12 u\log u\), while the Euler-product exponent above is
\[
  O(\log\log X+\sqrt u/\log u)=o(u\log u).
\]
Therefore
\[
  \Psi(2X,Y_X)\le 2X\exp(-\Omega(u\log u)),
\]
which beats every fixed power of \(\log X\) because \(\rho>1\).
\end{proof}

\begin{lemma}[Actual blocks and final suffixes]\label{lem:monotonicity}
For any strictly increasing accepted set:
\begin{enumerate}[label=(\alph*)]
\item for fixed accepted endpoint \(x\) and fixed length \(\ell\), there is
at most one actual consecutive accepted block of length \(\ell\) ending at
\(x\);
\item for fixed \(r\ge0\), the function
\[
  F_r(n)=n\prod_{\substack{a<n\\a\ {\rm accepted}\\a\ {\rm among\ last}\ r}}
  a
\]
is strictly increasing on candidates \(n\) for which the last \(r\) accepted
predecessors exist.
\end{enumerate}
\end{lemma}

\begin{proof}
Part (a) is immediate from strict increase.  If \(n_1<n_2\), then the ordered
\(r\)-tuple of last accepted predecessors before \(n_1\) is coordinatewise
\(\le\) the corresponding tuple before \(n_2\).  The predecessor product is
therefore nondecreasing, and multiplication by \(n_1<n_2\) gives
\(F_r(n_1)<F_r(n_2)\).
\end{proof}

\section{Private Barriers}\label{sec:barriers}

We turn to the construction of the deterministic set \(D\).  The barriers are
built only here; everything afterwards rests solely on the three properties
recorded in Proposition~\ref{prop:barriers}.  Put
\[
  \beta=6,\qquad L_X=\floor{(\log X)^\beta},\qquad
  L_0(X)=(\log(3X))^{2.1}.
\]
For each sufficiently large dyadic shell \((X,2X]\), partition the shell into
full consecutive intervals of length \(L_X\); boundary remainders are put in
\(D\).  In a full block \(I=[U,U+L_X)\), an integer start \(y\) is
\emph{admissible} if \(U\le y\) and \(y+L_0(X)<U+L_X\).

\begin{proposition}[Finite-horizon private barriers]\label{prop:barriers}
There is a deterministic set \(D\subset\N\) of density zero and a set of
selected semiprime barriers \(b=p_1p_2\) such that:
\begin{enumerate}[label=(\roman*)]
\item every selected barrier has
\[
  (\log b)^{1.08}<p_1\le(\log b)^{1.1},
\]
and its large prime \(p_2\) is a globally unique private label;
\item every selected barrier lies outside \(D\) and is accepted by the
modified greedy construction;
\item for every sufficiently large \(N\), put
\[
  F_Y=Y^{\floor{\sqrt{\log\log Y}}},
\]
and let \(X_0=X_0(N)\) be the dyadic scale with
\(F_{X_0}<N\le F_{2X_0}\).  Outside \(D\), a finite initial segment, and the
\(o(N)\) inverse-horizon smooth family whose chosen witness has old-left
maximum \(x\le5X_0\), every nonforced nonempty rejection \(n\le N\), with its
chosen overlap-cancelled witness, has a previous selected barrier \(q_B(n)<n\)
such that
\[
  n-q_B(n)=O((\log n)^\beta),
\]
the final suffix of the witness starts after \(q_B(n)\), and the old-left
block in the witness has value span \(O((\log n)^\beta)\).
\end{enumerate}
\end{proposition}

Only the properties listed in Proposition~\ref{prop:barriers} are used in the
later localization and counting arguments.

\begin{proof}
We start above a large fixed dyadic shell, put \(1\in D\), keep \(2\notin D\),
and absorb all finite initial choices into \(O(1)\).

\smallskip
\noindent\emph{Good blocks and label supply.}
Call a full \(L_X\)-block bad if fewer than half of its admissible starts are
good in the sense of Input~\ref{inp:MT}.  Since the admissible-start sets of
distinct full blocks are disjoint and each full block has \((1+o(1))L_X\)
admissible starts, Input~\ref{inp:MT}, applied up to \(2X+O(L_X)\), implies
that bad blocks in \((X,2X]\) have total length \(O(X/(\log X)^\delta)\).
Put all bad blocks in \(D\).
For every admissible start in the shell, \(y<3X\) for large \(X\), so
\((\log y)^{2.1}\le L_0(X)\); hence the MT interval attached to \(y\) remains
inside its long block.

In any remaining block, the good starts contribute
\(\gg L_X(\log X)^{1.1}\) start-semiprime incidences.  A fixed semiprime
\(b\) can be counted only for starts \(y\) with \(b-L_0(X)<y<b\), hence for at
most \(L_0(X)+O(1)\) starts.  Thus the block contains at least
\[
  \gg \frac{L_X(\log X)^{1.1}}{L_0(X)}
  =(\log X)^{\beta-1-o(1)}
\]
distinct semiprimes.  A fixed large prime label \(p_2\) gives at most one such
semiprime in one long block: two products \(p_1p_2\), \(p'_1p_2\) differ by
at least \(p_2\), while for these semiprimes
\(p_2=b/p_1\gg X/(\log X)^{1.1}\gg L_X\).  Hence each good block has
\((\log X)^{\beta-1-o(1)}\) distinct candidate labels.
For every counted \(b=p_1p_2\), since \(b\in(y,y+L_0(X)]\), we have
\(\log b=\log y+O(L_0(X)/X)=(1+o(1))\log X\).  The upper bound below is
immediate, and the lower bound follows from the exponent gap \(1.09>1.08\);
after increasing the initial cutoff,
\[
  (\log b)^{1.08}<p_1\le(\log b)^{1.1}.
\]

\smallskip
\noindent\emph{Greedy block-label matching.}
Order good blocks by left endpoint and candidate labels in each block by
increasing \(p_2\).  Scan blocks in order and match a block to the first
candidate label which was not selected in an earlier shell and has not already
been matched in the current shell.  If no such label exists, leave the block
unmatched.

A label selected in an earlier shell \((Y,2Y]\) can also be a candidate in
\((X,2X]\) only for polylogarithmically comparable \(Y\).  Comparing the two
size ranges
\[
  \frac{X}{(\log X)^{1.1}}\ll p_2\ll \frac{X}{(\log X)^{1.08}},
  \qquad
  \frac{Y}{(\log Y)^{1.1}}\ll p_2\ll \frac{Y}{(\log Y)^{1.08}},
\]
first gives \(Y=X^{1+o(1)}\), and then the safer window
\[
  cX(\log X)^{-0.05}\le Y\le C X(\log X)^{0.05}.
\]
Thus the number of earlier selected labels relevant to the current shell is
\[
  O\!\left(\frac{X}{L_X}(\log X)^{0.06}\right),
\]
where the factor \((\log X)^{0.01}\) absorbs the \(O(\log\log X)\) dyadic
shells in the comparable range.  Current-shell matched labels are absorbed in
the same bound.  Also a fixed label is adjacent to at most
\((\log X)^{1.1+o(1)}\) current blocks: \(p_1\) has only that many possible
prime values, and each product \(p_1p_2\) lies in \(O(1)\) long blocks.  If
\(U\) good blocks remain unmatched, then every candidate label of each
unmatched block is either an earlier selected label or a label already matched
in the current shell.  Counting these incidences gives
\[
  U(\log X)^{\beta-1-o(1)}
  \ll \frac{X}{L_X}(\log X)^{0.06}(\log X)^{1.1+o(1)}.
\]
Hence
\[
  U=O\!\left(\frac{X}{L_X}
  (\log X)^{-(\beta-2.16)+o(1)}\right)=o(X/L_X).
\]
Put all unmatched good blocks in \(D\).  In each matched block select the
matched semiprime as its barrier.

\smallskip
\noindent\emph{Prefix deletions.}
After every maximal deleted run of long blocks, and after each shell boundary,
look at the first subsequent matched block and delete its nonbarrier prefix
before the selected barrier.  The number of bad blocks is
\(O(X/(L_X(\log X)^\delta))\), and the number of unmatched blocks is
\[
  O\!\left(\frac{X}{L_X}
  (\log X)^{-(\beta-2.16)+o(1)}\right).
\]
Multiplying by the prefix length \(L_X\), the block and prefix deletions in
\((X,2X]\) have size
\[
  O(X/(\log X)^\delta)
  +O\!\left(X(\log X)^{-(\beta-2.16)+o(1)}\right)
  +O(L_X)=o(X).
\]

\smallskip
\noindent\emph{Finite-horizon multiples.}
For a selected barrier in shell \(X\), let
\[
  C_X=\floor{\sqrt{\log\log X}},\qquad F_X=X^{C_X}.
\]
For its private label \(p_2\), put into \(D\) every multiple
\(m\ne b\) of \(p_2\) with \(m\le F_X\).  These horizon deletions have density
zero.  A shell \((X,2X]\) contributes at most
\[
  O\!\left(\frac{\min(N,F_X)}{(\log X)^{\beta-1.1}}\right)
\]
to \(D\cap[1,N]\), since there are \(O(X/L_X)\) selected labels and each
label is \(\gg X/(\log X)^{1.1}\).  Put \(s=\beta-1.1>1\).

For old shells with \(F_X<N\), the dyadic sum is dominated by the last such
shell because \(C_X\) is nondecreasing and
\[
  F_{2X}/F_X\ge 2^{\frac12\sqrt{\log\log X}}
\]
eventually.  This gives \(O(N/(\log X_0)^s)\), where \(X_0\) is the largest
dyadic shell with \(F_{X_0}<N\).  Recent shells \(X_0<X\le N\) contribute
\[
  N\sum_{X_0<X\le N}\frac1{(\log X)^s}
  =O\!\left(\frac{N}{(\log X_0)^{s-1}}\right)=o(N),
\]
because \(F_{X_0}<N\le F_{2X_0}\) implies
\(\log X_0\ge(1+o(1))\log N/\sqrt{\log\log N}\).  Future shells \(X>N\) can
affect \([1,N]\) only if \(X/(\log X)^{1.1}\le N\), hence
\(X\le N(\log N)^2\) for large \(N\); their total contribution is
\[
  O\!\left(\frac{N\log\log N}{(\log N)^s}\right)=o(N).
\]
Finally, the other deterministic deletions have dyadic total \(o(N)\).  The
boundary remainders contribute \(O((\log N)^{\beta+1})\), and the MT-bad blocks
contribute \(O(N/(\log N)^\delta)\).  The unmatched-block and prefix
contribution is \(O(X(\log X)^{-(\beta-2.16)+o(1)})\) in a shell; since
\(\beta=6\), its dyadic total is also \(o(N)\).  This proves
\(\abs{D\cap[1,N]}=o(N)\).

\smallskip
\noindent\emph{Barrier survival and acceptance.}
A selected barrier is in a matched block, is not in a deleted prefix, and is
protected from its own label deletion.  Nor can any other selected label
delete it: if a selected label \(p\) from shell \(Y\) divides \(b=p_1p_2\) and
\(b\le F_Y\), then \(p=p_1\) or \(p=p_2\).  The case \(p=p_2\) contradicts
global label uniqueness.  The case \(p=p_1\) is impossible for large \(Y\),
because
\[
  p\ge \frac{Y}{(\log Y)^{1.1}},
  \qquad
  p_1\le(\log b)^{1.1}\le(\log F_Y)^{1.1}
  =(C_Y\log Y)^{1.1}
\]
and the last quantity is
\[
  o\!\left(\frac{Y}{(\log Y)^{1.1}}\right).
\]

When \(b\) is scanned, every earlier multiple of \(p_2\) has been deleted,
since \(b\le2X<F_X\).  Thus no old accepted product is divisible by \(p_2\):
if \(p_2\) divided such a product, it would divide one old accepted factor,
which would be an earlier deleted multiple of \(p_2\).  Every new product
containing \(b\), on the other hand, is divisible by \(p_2\).  Hence no
old-new collision can reject \(b\).  Old-old collisions are excluded by
Lemma~\ref{lem:greedy-unique}.  Finally, two new products containing \(b\)
would, after cancelling \(b\), give equality between two final suffix products
of the old accepted prefix; distinct final suffixes are strictly ordered by
inclusion and each proper extension multiplies by an accepted term at least
\(2\).  Hence no new-new collision can reject \(b\).
Therefore each selected barrier is accepted.

\smallskip
\noindent\emph{Blocking.}
Consider a surviving overlap-cancelled witness
\[
  \prod_{\text{old left}} d_i
  =n\prod_{\text{final suffix}} d_i
\]
with the old-left block strictly before the final suffix.  Suppose the old
left block contains a selected barrier \(b=p_1p_2\) from shell \(X\), and
\(n\le N\le F_X\).  Then \(p_2\) divides the left product, hence divides the
right product.  The corresponding terminal factor is either \(n\) or a suffix
factor.  It is not \(b\), because the right side lies strictly to the right of
the old block.  It is not any selected barrier \(c=r_1r_2\le F_X\): if
\(p_2=r_2\), global label uniqueness is violated; if \(p_2=r_1\), then
\[
  p_2\le(\log c)^{1.1}\le(\log F_X)^{1.1}
  =o(X/(\log X)^{1.1}),
\]
contradicting \(p_2\gg X/(\log X)^{1.1}\).  Thus the terminal factor is a
nonbarrier multiple of \(p_2\) below \(F_X\), hence belongs to \(D\), which is
impossible for either \(n\) or an accepted suffix factor.

\smallskip
\noindent\emph{Previous barriers and localization.}
Every sufficiently large nondeleted integer \(y\) which is not itself a
selected barrier has a previous selected barrier \(q_B(y)<y\) with
\[
  y-q_B(y)=O((\log y)^\beta).
\]
If \(y\) lies after the selected barrier in its matched block, use that
barrier.  If \(y\) lies before it, then \(y\) is in the nonbarrier prefix of
that matched block.  Since \(y\notin D\), this block cannot be the first
matched block after a deleted run or after the shell boundary, because those
prefixes were deleted.  The preceding full block is therefore undeleted and
matched, and its barrier is below \(y\) within two \(L_X\)-blocks.

Let \(n\) be a surviving nonempty rejection.  It is not itself a selected
barrier, since selected barriers are accepted.  The final suffix in its
chosen witness contains no selected barrier: if it contained \(b=p_1p_2\),
then every old-left factor would be \(<b\le2X_b<F_{X_b}\), where \(X_b\) is
the shell of \(b\).  An old-left factor divisible by \(p_2\) could not be a
selected barrier, by global label uniqueness and the small-factor bound
already used above; it would therefore be a deleted nonbarrier multiple of
\(p_2\).  Thus the left product would not be divisible by \(p_2\), while the
right product would be, a contradiction.  If the final suffix began at or
before \(q_B(n)\), it would contain the accepted selected barrier \(q_B(n)\).
Consequently the suffix starts after \(q_B(n)\), and its length is
\(O((\log n)^\beta)\).

Finally choose, for large \(N\), the dyadic \(X_0=X_0(N)\) with
\[
  F_{X_0}<N\le F_{2X_0}.
\]
If the old-left maximum \(x\) in the chosen witness satisfies \(x\le5X_0\),
then \(n\) is \(5X_0\)-smooth, because \(n\) divides the old-left product.
Moreover
\[
  \frac{\log N}{\log(5X_0)}
  =(1+o(1))\sqrt{\log\log N}.
\]
Input~\ref{inp:dickman} shows that these endpoints are \(o(N)\).  For all
remaining endpoints, \(x>5X_0\).  If \(x\) is a selected barrier, its horizon
is at least \(N\): the dyadic scale of \(x\) is above \(2X_0\), and
\(X\mapsto F_X=X^{\floor{\sqrt{\log\log X}}}\) is eventually increasing on
dyadic scales, so this horizon is at least \(F_{2X_0}\ge N\).  The blocking
argument forbids survival.  If \(x\) is not a selected barrier, the previous
barrier below \(x\) lies within \(O((\log x)^\beta)\) and, for large \(N\),
above \(2X_0\), since \(x\le N\le F_{2X_0}\) makes every fixed power of
\(\log x\) equal to \(o(X_0)\).  The same monotonicity gives horizon at least
\(N\).  The old-left block cannot contain that previous barrier, so it must
start after it.  Thus the old-left span is \(O((\log n)^\beta)\).
\end{proof}

\section{Localized Nonempty Witnesses}\label{sec:localization}

Fix the deletion set and selected barriers from Proposition~\ref{prop:barriers},
and for a nonempty rejection \(n\) use the deterministic overlap-cancelled
witness of Section~\ref{sec:greedy}.  What remains is to turn the qualitative
barrier information into a concrete numerical product identity, one in which
every local shift is bounded by a fixed power of \(\log X\).

\begin{proposition}[Localized witness]\label{prop:localized}
There is a fixed exponent \(B>20\) such that, in every sufficiently large
dyadic shell \(X<n\le2X\), all but \(o(X)\) nonempty rejection endpoints have
a chosen witness
\[
  \prod_{i=1}^{\ell}(x-a_i)=\prod_{j=1}^{b}(q+c_j),
  \qquad q=q_B(n),\quad n=q+h,
\]
with the following properties:
\begin{align*}
&2\le h\le(\log X)^7,\qquad
0\le a_i<T_X,\qquad 1\le c_j\le h,\qquad T_X=(\log X)^B,\\
&b=r+1,\qquad \ell=b+s,\qquad s\ge1,\qquad r\le(\log X)^7,\\
&\ell\le \ceil{((\log X)^7+1)\log_2(2X)}+2,\\
&b\ge c\,\frac{\log X}{\log\log X}\quad\text{for a fixed }c>0,\\
&x=q^{b/\ell}+O(T_X)\quad\text{with an absolute implicit constant},\\
&(\log X)^{-B}\le 1-\frac b\ell\le B\frac{\log\log X}{\log X},\\
&x<q<2X.
\end{align*}
The left factors \(x-a_i\) are distinct accepted terms, the terminal values
\(q+c_j\) are distinct, and they are exactly \(n=q+h\) together with the
actual final length-\(r\) suffix before \(n\).
\end{proposition}

\begin{proof}
The witness split gives an equality between an old block and \(n\) times a
final suffix.  The old block cannot be wholly contained in that suffix, since
then its product would be at most the suffix product, whereas the right side is
\(n\) times the suffix product with \(n>1\).  If it overlaps the suffix,
cancel the common contiguous overlap.  Since the suffix is final, these are
the only possibilities besides the old block already lying strictly before
the suffix.  If cancellation exhausted the suffix, then \(n\) would be an old
block product and hence an empty-suffix rejection.  Thus the chosen witness has
the form
\[
  \prod_{t=u}^{v} d_t
  = n\prod_{t=c}^{m} d_t,\qquad v<c\le m.
\]
Let \(x=d_v\), let \(w=d_u\), and let \(r=m-c+1\).  Proposition~\ref{prop:barriers},
applied with ambient bound \(2X\), removes \(o(X)\) deterministic,
inverse-horizon smooth, and finite-initial endpoints and gives
\begin{align*}
  2\le h=n-q_B(n)&\le(\log X)^7,\\
  r&\le(\log X)^7,\\
  x-w&\le(\log X)^7,\\
  x&>(\log X)^7.
\end{align*}
After this removal, the remaining endpoints have \(x>5X_0(2X)\).  This
inverse-horizon scale satisfies
\(\log X_0(2X)=(1+o(1))\log X/\sqrt{\log\log X}\), so
\(X_0(2X)>(\log X)^C\) for every fixed \(C\), once \(X\) is large; in
particular \(x>(\log X)^7\).
The lower bound \(h\ge2\) follows because the nonempty suffix starts after
\(q_B(n)\) and contains an accepted term \(<n\).

The terminal side has \(r+1\) factors, each at most \(2X\), while every
old-left factor is an accepted integer at least \(2\).  Hence
\[
  2^\ell\le (2X)^{r+1},
\]
which gives the stated upper bound for \(\ell\).  Also \(\ell\ge r+2\):
all terminal factors are larger than all old-left factors, so if
\(\ell\le r+1\), the terminal product would strictly exceed the old product.

Two elementary tails are negligible.  First suppose
\[
  \frac{\ell}{r+1}\ge 1+60\frac{\log\log X}{\log X}.
\]
Since the old-left span is at most \((\log X)^7\),
\[
  (x-(\log X)^7)^\ell\le(2X)^{r+1},
\]
so \(x\le(\log X)^7+X/(\log X)^{40}\le X/(\log X)^{39}\) for large \(X\).
For fixed \(x,\ell,r\), the
old-left block and its product are fixed, and the right side
\[
  n\cdot\text{(last \(r\) accepted predecessors before \(n\))}
\]
is strictly increasing by Lemma~\ref{lem:monotonicity}.  Thus there is at
most one endpoint for each fixed \(x,\ell,r\).  Here \(x\) has
\(O(X/(\log X)^{39})\) possible values, while \(\ell\) and \(r\) have only
fixed-polylogarithmically many possible values; the total is therefore
\(O(X/(\log X)^{20})=o(X)\).  Second, if
\[
  r\le \floor{\frac{\log X}{60\log\log X}},
\]
then the same argument gives
\[
  x\le(\log X)^7+(2X)^{(r+1)/(r+2)}\le \frac{3X}{(\log X)^{50}},
\]
and the same fixed-product and monotonicity argument, now with only
polylogarithmically many \((\ell,r)\) and \(O(X/(\log X)^{50})\) possible
values of \(x\), gives \(o(X)\) endpoints.

For the remaining endpoints, with \(b=r+1\),
\[
  b\ge c\frac{\log X}{\log\log X},\qquad
  \ell=b+s,\qquad
  1\le s<60b\frac{\log\log X}{\log X}.
\]
Choose fixed \(B>20\) so large that \(T_X=(\log X)^B\) dominates all fixed
polylogarithmic constants in the preceding estimates.  The witness becomes
\[
  \prod_{i=1}^{\ell}(x-a_i)=\prod_{j=1}^{b}(q+c_j),
\]
where \(0\le a_i<T_X\), \(1\le c_j\le h\), and \(n=q+h\) is included among
the terminal values.

The terminal product lies between \(q^b\) and \((q+h)^b\), and the old product
lies between \((x-(\log X)^7)^\ell\) and \(x^\ell\).  Thus
\[
  q^{b/\ell}\le x\le(\log X)^7+(q+h)^{b/\ell}.
\]
Since \(0<b/\ell<1\), the function \(t^{b/\ell}\) is increasing and concave on
\([q,q+h]\), with derivative at most \(1\).  Hence
\((q+h)^{b/\ell}-q^{b/\ell}\le h\), and therefore
\[
  0\le x-q^{b/\ell}\le(\log X)^7+h\le T_X
\]
for large \(X\), after enlarging \(B\) if necessary.
The lower bound
\[
  1-\frac b\ell\ge(\log X)^{-B}
\]
follows from \(s\ge1\) and \(\ell\le(\log X)^B\); the upper bound follows
from \(s<60b\log\log X/\log X\), after enlarging \(B\).

Finally let \(\theta=b/\ell\) and \(\Delta=1-\theta\ge(\log X)^{-B}\).
Since \(q\ge X-(\log X)^7\), the function
\(\Delta\mapsto 1-e^{-\Delta\log q}\) is increasing, and
\((\log X)^{-B}\log q=o(1)\), we have
\[
  q-q^\theta=q(1-e^{-\Delta\log q})
  \gg \frac{X\log X}{(\log X)^B}\gg T_X.
\]
Together with \(0\le x-q^\theta\le T_X\), this gives \(x<q\).  The inequality
\(q<2X\) is immediate from \(q<n\le2X\).
\end{proof}

\section{Terminal Support}\label{sec:terminal}

A single product identity does not by itself bound the number of endpoints;
what it yields is arithmetic information.  Any large prime dividing the
left-hand side of the localized identity must also divide one of the terminal
factors on the right.  The two lemmas of this section render this observation
quantitative, extracting from each identity a family of independent
large-prime congruences, which serve as the input to the counting argument of
Section~\ref{sec:slab}.

\begin{lemma}[Smooth accepted left factors]\label{lem:smooth-left}
Let
\[
  Y_X=\exp\!\left(\frac{\log X}{(\log\log X)^2}\right).
\]
Among the localized endpoints of
Proposition~\ref{prop:localized}, those for which some accepted left factor
\(x-a_i\) has \(\pplus(x-a_i)\le Y_X\) form an \(o(X)\) family in the shell
\((X,2X]\).
\end{lemma}

\begin{proof}
By Lemma~\ref{lem:smooth-decay}, the number of \(Y_X\)-smooth integers
\(\le2X\) is
\[
  o\!\left(\frac{X}{(\log X)^M}\right)
\]
for every fixed \(M\).  Given such an
integer \(m=x-a_i\), there are at most \(T_X\) choices for \(x\).  For fixed
\(x\) and \(\ell\) there is at most one actual left block, and for each
fixed suffix length \(r\) at most one endpoint.  For the fixed left block has a
fixed product \(P\), while any endpoint counted with that block and this \(r\)
must satisfy
\[
  P=n\prod_{\substack{a<n\\a\ {\rm accepted}\\a\ {\rm among\ last}\ r}}a,
\]
and the right side is strictly increasing in \(n\) by
Lemma~\ref{lem:monotonicity}.
These at most \(T_X\) choices of \(x\) already account for the possible position of
the smooth factor inside the left block.  The remaining choices of \(\ell\)
and \(r\) are at most \(O((\log X)^8)\) and \(O((\log X)^7)\), respectively.
Thus the endpoint count is
\[
  o\!\left(\frac{X}{(\log X)^M}\right)O((\log X)^{B+15}),
\]
and taking \(M\) larger than this fixed loss gives \(o(X)\).
\end{proof}

\begin{lemma}[Terminal-support extraction]\label{lem:terminal-support}
Consider a localized endpoint from Proposition~\ref{prop:localized}, not
counted in Lemma~\ref{lem:smooth-left}, and take \(X\) large enough that
\((\log X)^7<T_X\) and \(T_X+(\log X)^7<Y_X\).
For every left factor choose
\[
  p_i=\pplus(x-a_i)>Y_X.
\]
Put
\[
  \Lambda_X=\ceil{\frac{\log(3X)}{\log Y_X}}+2.
\]
Then the endpoint supplies at least \(\floor{\ell/\Lambda_X^2}\) distinct
local slab points \((a_i,r_i)\), \(r_i=a_i+c_i\), with distinct \(r_i\)'s and
carrying distinct prime labels \(p_i>Y_X\) such that
\[
  p_i\mid x-a_i,\qquad p_i\mid D_0+r_i,\qquad D_0=q-x>0.
\]
\end{lemma}

\begin{proof}
Since \(p_i\) divides the left product, it divides one terminal value
\(q+c_i\), \(1\le c_i\le h\).  This terminal value is unique, because two
terminal values differ by less than \(h<Y_X<p_i\).  A label \(p_i\) appears
on at most one left factor and at most one terminal value, since left factors
differ by less than \(T_X<Y_X\) and terminal values differ by at most \(h<Y_X\).

A fixed terminal value \(\le n\le2X<3X\) is incident to at most
\(\Lambda_X\) chosen labels, since these labels are distinct primes and the
value cannot have more than
\(\log(3X)/\log Y_X+1\) distinct prime factors exceeding \(Y_X\).  Therefore
the incidence graph from \(\ell\) left vertices to terminal vertices contains
a matching of size at least \(\floor{\ell/\Lambda_X}\): the \(\ell\) left
edges occupy at least \(\ceil{\ell/\Lambda_X}\) terminal vertices, and choosing
one incident edge over each occupied terminal gives a matching, since each
left vertex has degree one.  The matched labels are distinct.

For a matched edge, \(p_i\mid x-a_i\) and \(p_i\mid q+c_i\), so
\[
  p_i\mid(q+c_i)-(x-a_i)=D_0+a_i+c_i.
\]
Set \(r_i=a_i+c_i\).  For a fixed value \(\rho\), all matched labels with
\(r_i=\rho\) divide the positive integer \(D_0+\rho\), where positivity follows
from \(D_0>0\) and \(\rho=a_i+c_i\ge1\).  Since
\(D_0+\rho\le q-x+T_X+h<3X\) for large \(X\), there
are at most \(\Lambda_X\) of them.  Selecting one representative for each
distinct \(r_i\) leaves at least \(\floor{\ell/\Lambda_X^2}\) slab points:
if \(M\ge\floor{\ell/\Lambda_X}\) matched edges remain and at most
\(\Lambda_X\) have the same \(r_i\), then
\(\ceil{M/\Lambda_X}\ge\floor{\ell/\Lambda_X^2}\): writing
\(\ell=a\Lambda_X^2+r\) with \(0\le r<\Lambda_X^2\), one has
\(\floor{\ell/\Lambda_X}\ge a\Lambda_X\), and hence
\(\ceil{\floor{\ell/\Lambda_X}/\Lambda_X}\ge a=\floor{\ell/\Lambda_X^2}\).
The retained points have distinct labels and the stated divisibilities.
\end{proof}

\section{The Shifted-CRT Slab Count}\label{sec:slab}

Once the terminal support is available, the remaining count is carried out by
a single proposition, whose proof rests on the determinant method.  One
interpolates a polynomial of low degree through the local points; the
congruences accumulated above then force it to vanish at the global point;
and a bound for the number of global points lying on a fixed curve of low
degree completes the count.  The argument divides into three parts: the
interpolation step, an estimate for the nonlinear components of the auxiliary
polynomial by means of Input~\ref{inp:CCDN}, and a separate treatment of its
line components, for which the determinant method is unavailable and the
curvature of \(q^\theta\) is used in its place.

\begin{lemma}[Auxiliary interpolation]\label{lem:interpolation}
Let \(Z\) be a set of \(K\ge1\) distinct integer points in a fixed
polylogarithmic box
\[
  \abs{A},\abs{R}\le(\log X)^{B_0},
\]
and assume \(K\log\log X=o(\log X)\).  With
\[
  m=\ceil{4\sqrt K},
\]
there is a nonzero nonconstant polynomial \(G(A,R)\in\Z[A,R]\), of total
degree at most \(m\) and height \(\exp(O(m\log\log X))\), vanishing on \(Z\).
It may be chosen primitive, with a fixed sign convention for its first
nonzero coefficient.
\end{lemma}

\begin{proof}
The number \(N_m\) of monomials \(A^uR^v\), \(u+v\le m\), is
\((m+1)(m+2)/2>2K\).  The \(K\times N_m\) evaluation matrix on \(Z\) has all
entries bounded by \(\exp(O(m\log\log X))\).  Compare coefficient vectors in
\([0,Q)^{N_m}\) under the \(K\) evaluation forms.  For
\[
  Q\ge (3N_mH_0)^{K/(N_m-K)}+1,\qquad H_0=\exp(O(m\log\log X)),
\]
two vectors have the same image; their nonzero difference gives an integer
kernel vector of height \(\exp(O(m\log\log X))\).  The corresponding
polynomial vanishes on \(Z\), and is not a nonzero constant because \(K\ge1\).
Dividing its coefficient vector by the gcd of the coefficients and fixing the
sign of the first nonzero coefficient preserves all these bounds and gives the
primitive normalized choice.
\end{proof}

\begin{lemma}[Line-factor height]\label{lem:line-height}
Let \(G(A,R)\in\Z[A,R]\) be primitive, nonzero, of total degree at most \(m\),
with
\[
  m^2=o(\log X),\qquad H(G)=\exp(o(\log X)).
\]
If
\[
  L(A,R)=\alpha A+\beta R+\gamma
\]
is a primitive integer line factor of \(G\) over \(\Q\), then
\[
  \max(\abs{\alpha},\abs{\beta},\abs{\gamma})=X^{o(1)}.
\]
\end{lemma}

\begin{proof}
Homogenize \(G\) as a degree-\(m\) form; the line becomes the
primitive homogeneous factor
\[
  L^h=\alpha A+\beta R+\gamma W.
\]
Permute variables so the largest coefficient is the coefficient \(c\) of the
homogenizing variable, and write the line as \(aU+bV+cW\).  Dehomogenizing at
\(W=1\) gives a polynomial \(P(U,V)\) of degree at most \(m\), height
\(\le H(G)\), divisible by \(aU+bV+c\).

If \(\max(|a|,|b|,|c|)=1\) there is nothing to prove.  Otherwise choose an
integer slope \(0\le t\le m+1\) such that \(a+bt\ne0\) and the cofactor in
\[
  P=(aU+bV+c)Q
\]
does not vanish identically on \((U,V)=(Z,tZ)\).  In this case
\((a,b)\ne(0,0)\).  Such a slope exists because at most \(m-1\) slopes kill
the nonzero cofactor and at most one slope makes \(a+bt=0\): writing
\(Q=\sum_j Q_j\) into homogeneous parts, if \(Q(Z,tZ)\) is identically zero,
then \(Q_j(1,t)=0\) for every \(j\), and one fixed nonzero \(Q_j\) has degree
at most \(m-1\).  Then
\[
  P_t(Z)=P(Z,tZ)
\]
is a nonzero one-variable integer polynomial of degree at most \(m\), divisible
by \((a+bt)Z+c\), and
\[
  \norm{P_t}_2\le \exp(O(m\log m))H(G).
\]
Mignotte's inequality, Input~\ref{inp:mignotte}, gives
\[
  \max(|a+bt|,|c|)\le 2^m\norm{P_t}_2
  \le\exp(O(m\log m))H(G).
\]
Since \(|c|\) was the largest line coefficient and
\(m\log m=O(m^2)=o(\log X)\), the result follows.
\end{proof}

\begin{lemma}[Nonabsolute factors]\label{lem:nonabsolute}
Let \(H(A,R)\in\Z[A,R]\) be primitive, squarefree, \(\Q\)-irreducible, not
absolutely irreducible, and of total degree \(d\ge2\).  Then
\[
  \#\{P\in\Q^2:H(P)=0\}\le d^2.
\]
\end{lemma}

\begin{proof}
Over \(\overline{\Q}\), the polynomial \(H\) is a product of at least two
distinct Galois-conjugate absolute factors.  A rational zero of \(H\) lies on
one absolute factor, hence by Galois conjugacy lies on every conjugate factor.
Choose two distinct factors.  They have no common component, so Bezout bounds
their common affine zero set by the product of their degrees, at most \(d^2\).
\end{proof}

\begin{proposition}[Shifted-CRT slab count]\label{prop:slab}
Fix constants \(A,B,B_1,\kappa>0\).  In a dyadic shell \(X<n\le2X\), suppose the
endpoints under consideration have integral localized data
\[
  n=q+h,\quad 1\le h\le(\log X)^A,\quad T_X=(\log X)^B,
\]
\[
  x=q^\theta+O(T_X),\quad D_0=q-x>0,\quad 0<\theta<1,
\]
where the implicit constant in \(O(T_X)\) is uniform over all endpoints with
the fixed parameter tuple.
\[
  (\log X)^{-B}\le 1-\theta\le B\frac{\log\log X}{\log X}.
\]
Assume that, for every fixed \(\eps>0\), the number of parameter tuples,
including \(h,\theta\), and \(K\), is \(O_\eps(X^\eps)\).  Let
\[
  Y_X=\exp(\log X/(\log\log X)^{\kappa}).
\]
Suppose each endpoint carries \(K\) distinct local slab points
\[
  z_i=(a_i,r_i),\qquad
  0\le a_i<T_X,\qquad |r_i|\le(\log X)^{B_1},
\]
with distinct prime labels \(p_i>Y_X\) satisfying
\[
  p_i\mid x-a_i,\qquad p_i\mid D_0+r_i.
\]
If
\[
  K\log\log X=o(\log X),\qquad
  \frac{K\log Y_X}{\sqrt K\,\log X}\to\infty,
\]
then the number of such endpoints in \((X,2X]\) is \(o(X)\).
\end{proposition}

\begin{proof}
The proof proceeds in the three parts outlined above.

\smallskip
\noindent\emph{Auxiliary curves.}
Fix one localized parameter tuple and one retained \(K\)-set \(Z\) of local
points.  The number of possible \(Z\)'s is \(X^{o(1)}\), since the local box is
polylogarithmic and
\[
  \binom{(\log X)^{O(1)}}{K}\le\exp(O(K\log\log X))=X^{o(1)}.
\]
Fix a constant \(C_G\), depending only on the local box exponents \(B\) and
\(B_1\), for
which Lemma~\ref{lem:interpolation} supplies a primitive nonzero polynomial
\(G(A,R)\) of degree at most \(m=\ceil{4\sqrt K}\), height at most
\(\exp(C_Gm\log\log X)\), and vanishing on \(Z\).  Among all such polynomials
with first nonzero coefficient positive, choose the lexicographically first
coefficient vector.  Thus \(G\) depends only on the fixed parameters and the
coordinate set \(Z\), not on the endpoint-specific prime labels.

For every endpoint and every retained point,
\[
  (x,-D_0)\equiv (a_i,r_i)\pmod {p_i}.
\]
Therefore each \(p_i\) divides \(G(x,-D_0)\), and the distinctness of the labels
gives
\[
  \prod_i p_i\mid G(x,-D_0).
\]
If \(G(x,-D_0)\ne0\), then, since \(G\) has \(O(m^2)=O(K)\) monomials and
\(|x|,|D_0|=O(X)\),
\[
  \log |G(x,-D_0)|=O(m\log X)=O(\sqrt K\,\log X).
\]
But
\[
  \log\prod_i p_i\ge K\log Y_X,
\]
and the ratio \(K\log Y_X/(\sqrt K\log X)\) tends to infinity.  This is
incompatible with the preceding bound for large \(X\).  Hence
\(G(x,-D_0)=0\).

\smallskip
\noindent\emph{Nonlinear components.}
Thus all endpoints lie on \(X^{o(1)}\) auxiliary curves \(G(A,R)=0\), each of
degree \(m=X^{o(1)}\).  Remove content and repeated factors, and then consider
the primitive squarefree factorization over \(\Q\).  This does not change the
zero set relevant to us: every point with \(G(x,-D_0)=0\) lies on at least one
primitive squarefree \(\Q\)-irreducible factor.  Primitive line factors are
deferred to the next paragraph.  Nonlinear absolutely irreducible factors of
degree \(d\ge2\) are integral affine curves.  For one such factor,
Input~\ref{inp:CCDN}, in the box of height \(O(X)\), gives
\[
  O(d^3X^{1/d}(\log X+d))\le X^{1/2+o(1)}
\]
points.  The constant in Input~\ref{inp:CCDN} depends only on the fixed ambient
dimension, and the squarefree factors of one auxiliary curve have total degree
at most \(m=X^{o(1)}\); summing this bound over the absolutely irreducible
nonlinear factors still gives \(X^{1/2+o(1)}\) points for that auxiliary curve.
If a \(\Q\)-irreducible nonlinear factor is not absolutely irreducible,
Lemma~\ref{lem:nonabsolute} gives \(O(d^2)\) rational points, so all
nonabsolute factors of one auxiliary curve contribute only \(X^{o(1)}\) points.
Multiplication by the \(X^{o(1)}\) curves still gives \(o(X)\) endpoints, since
the fixed parameter tuple includes \(h\), and \((x,-D_0)\) determines
\(q=x+D_0\) and \(n=q+h\).

\smallskip
\noindent\emph{Line components.}
It remains to count primitive line factors
\[
  \alpha A+\beta R+\gamma=0.
\]
Every such factor divides the original primitive auxiliary polynomial \(G\)
over \(\Q\), even after passing to the squarefree radical, so
Lemma~\ref{lem:line-height} applies to the same \(G\).
By Lemma~\ref{lem:line-height},
\[
  M_X:=\max(|\alpha|,|\beta|,|\gamma|)=X^{o(1)}
\]
uniformly for all such factors.  At the global point \(R=-D_0=x-q\),
\[
  (\alpha+\beta)x-\beta q+\gamma=0.              \tag{1}\label{eq:line-global}
\]
Also \(q\ge X/2\).  Since \(x=q^\theta+O(T_X)\) and
\[
  q^\theta\ge q/(\log X)^{O(B)}\gg T_X,
\]
we have
\[
  x\gg q/(\log X)^{O(B)}.
\]
For large \(X\), neither \(\beta\) nor \(\alpha+\beta\) can be zero in
\eqref{eq:line-global}, otherwise \(x\) or \(q\) would be \(X^{o(1)}\).

Define
\[
  F(q)=(\alpha+\beta)q^\theta-\beta q+\gamma.
\]
Since \(x=q^\theta+O(T_X)\), equation \eqref{eq:line-global} gives, for a
fixed constant \(C_1\),
\[
  |F(q)|\le C_1M_XT_X.                           \tag{2}\label{eq:line-sublevel}
\]
On the sublevel set \eqref{eq:line-sublevel}, we have
\[
  |F'(q)|\ge c|\beta|(1-\theta).                 \tag{3}\label{eq:line-derivative}
\]
To verify this, suppose \(F'(q)=\delta\) with
\(|\delta|<|\beta|(1-\theta)/4\); then
\[
  (\alpha+\beta)q^\theta=\frac{(\beta+\delta)q}{\theta}
\]
and hence
\[
  F(q)=\frac{\beta q(1-\theta)}{\theta}
       +\frac{\delta q}{\theta}+\gamma.
\]
Using \(|\beta|\ge1\), \(q\gg X\),
\((1-\theta)\ge(\log X)^{-B}\), and \(\theta\ge1/2\), the first term in this
last display has magnitude \(\gg X/(\log X)^B\).  The \(\delta\)-term has at
most one quarter of that magnitude, and \(|\gamma|\le M_X=X^{o(1)}\) is
negligible in comparison.  Hence \(|F(q)|\gg X/(\log X)^B\), which is much
larger than \(C_1M_XT_X\).
This contradicts \eqref{eq:line-sublevel}, proving
\eqref{eq:line-derivative}.

The derivative \(F'\) is monotone on \(q>0\), since \(0<\theta<1\) and
\(\alpha+\beta\ne0\), so \(F'\) has at most one zero.  Thus \(F\) is monotone
on at most two intervals, and the sublevel set
\(|F(q)|\le C_1M_XT_X\) has at most two interval components.  Each component
has length
\[
  O\!\left(\frac{C_1M_XT_X}{|\beta|(1-\theta)}\right)=X^{o(1)}
\]
by the mean-value theorem and \eqref{eq:line-derivative}.  Hence each line
factor contributes \(X^{o(1)}\) integer values of \(q\), and, since \(h\) is
fixed by the parameter tuple, \(X^{o(1)}\) endpoints \(n=q+h\).  There are only
\(X^{o(1)}\) curves and line factors.  Thus all line components contribute
\(o(X)\).  Combining the nonlinear and line estimates proves the proposition.
\end{proof}

\section{Counting Nonempty Rejections}\label{sec:nonempty}

The pieces are now in place.  Every exceptional family has already been shown
to have shell size \(o(X)\), and each endpoint that survives them carries
enough terminal support to be counted by Proposition~\ref{prop:slab}.

\begin{proposition}\label{prop:nonempty}
For the modified greedy construction using the deterministic deletion set from
Proposition~\ref{prop:barriers}, the number of nonempty rejections \(n\le N\)
is \(o(N)\).
\end{proposition}

\begin{proof}
Work in a dyadic shell \(X<n\le2X\).  Remove the deterministic deletions,
empty-suffix endpoints, and the exceptional families from
Propositions~\ref{prop:barriers} and~\ref{prop:localized}.  These contribute
\(o(X)\) in the shell: the deterministic deletions are density zero, the
empty-suffix endpoints are controlled by Corollary~\ref{cor:empty-suffix} with
ambient bound \(2X\), and the localization exceptions were proved either as
shell \(o(X)\) families or as global \(o(N)\) families applied with \(N=2X\).
These estimates were established independently of the present proposition.

For every remaining endpoint, Proposition~\ref{prop:localized} applies.
Discard the \(o(X)\) smooth-left family of Lemma~\ref{lem:smooth-left}.  Then
every left factor has a chosen label \(p_i=\pplus(x-a_i)>Y_X\), where
\[
  Y_X=\exp(\log X/(\log\log X)^2).
\]
By Lemma~\ref{lem:terminal-support}, the endpoint supplies at least
\(\floor{\ell/\Lambda_X^2}\) slab points, where
\[
  \Lambda_X=\ceil{\frac{\log(3X)}{\log Y_X}}+2
  =(\log\log X)^2+O(1).
\]
Since Proposition~\ref{prop:localized} gives
\[
  \ell\ge b\ge c\frac{\log X}{\log\log X},
\]
we may retain the lexicographically first
\[
  K_{\rm slab}=
  \floor{\frac{c\log X}{100(\log\log X)^5}}
\]
of these extracted slab points.  The factor \(100\) leaves room for the floor
loss: for large \(X\), \(\Lambda_X\le2(\log\log X)^2\), so
\[
  \frac{\ell}{\Lambda_X^2}\ge
  \frac{c\log X}{4(\log\log X)^5},
\]
and hence
\[
  \frac{\ell}{\Lambda_X^2}-1
  \ge \frac{c\log X}{100(\log\log X)^5}
  \ge K_{\rm slab}.
\]
Thus \(\floor{\ell/\Lambda_X^2}\ge K_{\rm slab}\).

The hypotheses of Proposition~\ref{prop:slab} hold with \(A=7\), \(\kappa=2\),
\(B_1=B+1\), \(K=K_{\rm slab}\), and \(\theta=b/\ell\).  Here
Proposition~\ref{prop:localized} supplies \(n=q+h\), \(h\le(\log X)^7\),
\(x=q^\theta+O(T_X)\) with a fixed implicit constant, \(D_0=q-x>0\),
\(0<\theta<1\), and the two required bounds on \(1-\theta\).
Lemma~\ref{lem:terminal-support} gives the distinct
slab points and distinct prime labels with the two divisibilities; retaining
the first \(K_{\rm slab}\) points preserves these properties.  Finally,
\[
  |r_i|=|a_i+c_i|\le T_X+(\log X)^7\le(\log X)^{B+1}
\]
for large \(X\).  The localized parameter tuple is determined by \(h\),
\((b,\ell)\), and \(K_{\rm slab}\).  Since \(h\le(\log X)^7\) and, by
Proposition~\ref{prop:localized},
\[
  1\le b,\ell\le \ceil{((\log X)^7+1)\log_2(2X)}+2=O((\log X)^8),
\]
this gives only
\(O((\log X)^{23})=X^{o(1)}\) possibilities.  Moreover
\[
  K_{\rm slab}\log\log X=o(\log X),
\]
and, since
\[
  K_{\rm slab}\ge \frac{c\log X}{200(\log\log X)^5}
\]
for all large \(X\),
\[
  \frac{K_{\rm slab}\log Y_X}{\sqrt{K_{\rm slab}}\log X}
  =\frac{\sqrt{K_{\rm slab}}}{(\log\log X)^2}\to\infty.
\]
Thus Proposition~\ref{prop:slab} counts the surviving endpoints by \(o(X)\).
Adding back the finitely many \(o(X)\) families removed above, every dyadic
shell contains \(o(X)\) nonempty rejections.

Finally sum dyadic shells.  If \(f_j\) is the ratio of the number of nonempty
rejections in \((2^j,2^{j+1}]\) to \(2^j\), then \(f_j\to0\).  For every
\(\eps>0\), choose \(J_0\) so that \(f_j\le\eps\) for \(j\ge J_0\).  The
finite initial shells contribute \(O(2^{J_0})\), while the dyadic tail up to
\(N\) contributes at most
\[
  \eps\sum_{2^j\le N}2^j\le2\eps N.
\]
Letting \(\eps\to0\) gives \(o(N)\).
\end{proof}

\section{Proof of the Main Theorem}\label{sec:main-proof}

\begin{proof}[Proof of Theorem~\ref{thm:main}]
Use the deterministic deletion set \(D\) from Proposition~\ref{prop:barriers}
and run the modified greedy construction.  By
Proposition~\ref{prop:barriers}, forced deletions contribute \(o(N)\) rejected
values.  By Corollary~\ref{cor:empty-suffix}, empty-suffix collision
rejections contribute \(o(N)\).  By Proposition~\ref{prop:nonempty},
nonempty collision rejections contribute \(o(N)\).

Every rejected \(n\) belongs to the union of these three families.  If
\(n\in D\), it is forced.  If \(n\notin D\), Lemma~\ref{lem:witness-split}
gives an old-new witness \(P=nS\).  If some witness has \(S=1\), \(n\) is
empty-suffix; otherwise the deterministic chosen witness has \(S>1\), and
\(n\) is counted by the nonempty family.  Disjointness is unnecessary.  Hence
the total rejected set has density zero, so the accepted set has density one.
Lemma~\ref{lem:greedy-unique} gives distinct consecutive-block products at
every finite stage.  The increasing enumeration of the accepted set proves the
theorem.
\end{proof}

\begin{thebibliography}{9}

\bibitem{ErGr80}
P. Erd\H{o}s and R. L. Graham,
\emph{Old and New Problems and Results in Combinatorial Number Theory},
Monographies de L'Enseignement Math\'ematique, no. 28,
Universit\'e de Gen\`eve, 1980.

\bibitem{Bloom421}
T. F. Bloom,
\emph{Erd\H{o}s Problem \#421},
\url{https://www.erdosproblems.com/421}.

\bibitem{MT}
K. Matom\"aki and J. Ter\"av\"ainen,
\emph{Almost primes in almost all short intervals II},
arXiv:2207.05038.

\bibitem{deBruijn}
N. G. de Bruijn,
\emph{On the number of positive integers not exceeding \(x\) and free of
prime factors greater than \(y\)},
Indag. Math. 13 (1951), 50--60.

\bibitem{Hildebrand}
A. Hildebrand,
\emph{Integers without large prime factors},
J. Th\'eor. Nombres Bordeaux 5 (1993), 411--484.

\bibitem{CCDN}
W. Castryck, R. Cluckers, P. Dittmann, and K. H. Nguyen,
\emph{The dimension growth conjecture, polynomial in the degree and without
logarithmic factors},
Algebra \& Number Theory 14 (2020), 2261--2294; arXiv:1904.13109.

\bibitem{Mignotte}
M. Mignotte,
\emph{An inequality about factors of polynomials},
Mathematics of Computation 28 (1974), 1153--1157.

\end{thebibliography}

\end{document}
